% Generated draft: requires mathematical and source review.
\documentclass[11pt]{article}
\usepackage[margin=25mm]{geometry}
\usepackage{fontspec}
\setmainfont{Latin Modern Roman}
\usepackage{amsmath,amssymb,mathtools,unicode-math}
\setmathfont{Latin Modern Math}
\usepackage{enumitem,needspace,xurl,hyperref}
\setlength{\parindent}{0pt}
\setlength{\parskip}{5pt}
\setlength{\emergencystretch}{3em}
\newcommand{\N}{\mathbb{N}}
\newcommand{\Z}{\mathbb{Z}}
\newcommand{\Q}{\mathbb{Q}}
\newcommand{\R}{\mathbb{R}}
\newcommand{\C}{\mathbb{C}}
\newcommand{\F}{\mathbb{F}}
\newcommand{\A}{\mathbb{A}}
\newcommand{\PP}{\mathbb{P}}
\newcommand{\E}{\mathbb{E}}
\newcommand{\eps}{\varepsilon}
\newcommand{\dd}{\,\mathrm{d}}
\begin{document}
\section*{UnsolvedMath ID: 800007}
\textbf{Source problem code:} catalogue entry
\subsection{Definitions and mathematical statement}
A semigroup is a nonempty set with an associative binary operation. A semigroup variety is a class of semigroups defined by identities. For a class \(\mathcal G\) of semigroups, let \(\operatorname{Var}(\mathcal G)\) denote the smallest semigroup variety containing \(\mathcal G\). A semigroup \(S\) is regular if
\[
\forall a\in S\;\exists b\in S\qquad aba=a.
\]
A semigroup is periodic if every element generates a finite cyclic subsemigroup; equivalently,
\[
\forall a\in S\;\exists m,n\in\N\qquad 1\leq m<n\quad\text{and}\quad a^m=a^n.
\]
A variety is periodic if every semigroup belonging to it is periodic.

\textbf{Problem (S. Kublanovsky).} Characterize the periodic semigroup varieties \(\mathcal V\) for which there exists a class \(\mathcal G\) of regular semigroups such that
\[
\mathcal V=\operatorname{Var}(\mathcal G).
\]
The requirement concerns the generating semigroups; it does not assert that every member of \(\mathcal V\) is regular. The source does not specify a required form of the characterization.

The supplied source states that the corresponding description for periodic varieties generated by completely \(0\)-simple semigroups was obtained in reference [Kubl4]. Here a semigroup \(S\) with zero is completely \(0\)-simple if \(S\neq\{0\}\), \(S^2\neq\{0\}\), its only ideals are \(\{0\}\) and \(S\), and it contains a primitive nonzero idempotent. A nonzero idempotent \(e\) is primitive if every nonzero idempotent \(f\) satisfying \(ef=fe=f\) equals \(e\). The bibliographic entry for [Kubl4] is absent from the supplied record, so this background attribution remains unverified.
\subsection{Short English statement}
Describe exactly which periodic varieties of semigroups can be generated by regular semigroups. The source attributes the problem to S. Kublanovsky and reports an existing description for the special case of completely \(0\)-simple generators, citing [Kubl4] without supplying its bibliographic entry.
\subsection{Sources}
\begin{itemize}
\item Local UnsolvedMath registry, permanent ID \(800007\), source problem code \texttt{catalogue entry}, title \emph{The Lyapin's Notebook --- Problem 7}. The registry labels the problem open and lists Bershadsky, Kublanovsky, and Mashevitzky under proposal attribution, with proposed year \(2026\). The statement itself attributes the problem to S. Kublanovsky.
\item Bershadsky, Kublanovsky, and Mashevitzky, \emph{The Lyapin's Notebook: Unsolved Problems in Semigroup Theory} (2026), Problem 7. Supplied source URL: \url{https://arxiv.org/abs/2604.04763}. The supplied provenance records access on \(2026\text{-}10\text{-}16\) and extraction from arXiv LaTeX; this access information is supplied provenance, not an independently verified visit.
\item Reference [Kubl4], cited in the supplied statement for the description of periodic varieties generated by completely \(0\)-simple semigroups. No author details, title, publication information, or URL for this reference are supplied.
\end{itemize}
\end{document}
